\documentclass[11pt]{article}

\usepackage[utf8]{inputenc}
\usepackage[T1]{fontenc}
\usepackage{amsmath,amssymb,amsfonts,amsthm}
\usepackage{mathtools}
\usepackage{geometry}
\geometry{margin=1in}
\usepackage[hidelinks]{hyperref}

\numberwithin{equation}{section}

\theoremstyle{plain}
\newtheorem{theorem}{Theorem}[section]
\newtheorem{conjecture}[theorem]{Conjecture}
\newtheorem{lemma}[theorem]{Lemma}
\newtheorem{proposition}[theorem]{Proposition}
\newtheorem{corollary}[theorem]{Corollary}

\theoremstyle{definition}
\newtheorem{definition}[theorem]{Definition}

\theoremstyle{remark}
\newtheorem{remark}[theorem]{Remark}

\DeclareMathOperator{\Mass}{Mass}
\DeclareMathOperator{\PD}{PD}

\newcommand{\R}{\mathbb{R}}
\newcommand{\C}{\mathbb{C}}
\newcommand{\Z}{\mathbb{Z}}
\newcommand{\Q}{\mathbb{Q}}
\newcommand{\N}{\mathbb{N}}
\newcommand{\eps}{\varepsilon}
\newcommand{\mhol}{m_{\mathrm{hol}}}
\newcommand{\Def}{\mathrm{Def}}
\newcommand{\Defcal}{\Def_{\mathrm{cal}}}

\title{\bfseries Calibrated Discretization and the Hodge Conjecture:\\
  The Proof Spine of a Four-Paper Series}

\author{
    Jonathan Washburn\thanks{Recognition Science, Recognition Physics Institute,
        Austin, Texas, USA. Email: \texttt{jon@recognitionphysics.org}.}
    \and
    Amir Rahnamai Barghi\thanks{Recognition Science, Recognition Physics Institute,
        Austin, Texas, USA. Email: \texttt{arahnamab@gmail.com}.}
}

\date{\today}

\begin{document}
\maketitle

\begin{abstract}
Let $X$ be a smooth complex projective manifold of complex dimension $n$.
Assuming the interface theorems established in three companion manuscripts on
local holomorphic manufacturing, coherent packet assembly, and exact-class
gluing/rounding~\cite{Washburn26I,Washburn26II,Washburn26III}, we derive the
final Hodge conclusion for rational $(p,p)$-classes.
The argument has three stages:
(1) a signed decomposition reduces any rational Hodge class to the cone-positive case;
(2) for cone-positive classes, the imported discretization machine produces
integral cycles with vanishing calibration defect in a fixed homology class;
(3) Federer--Fleming compactness yields a calibrated holomorphic limit, which is
algebraic by Harvey--Lawson/King and Chow/GAGA.

This paper therefore contains only the proof spine of the series:
the closure chain from the interface theorems to the final Hodge conclusion.
The local holomorphic manufacturing, coherent discretization, and exact-class
gluing machinery are not reproved here; they are used only through the explicit
interface statements of Section~\ref{sec:interface}.
\end{abstract}


%%% ============================================================
\section{Introduction}\label{sec:intro}
%%% ============================================================

We formulate the Hodge problem for a fixed rational $(p,p)$ class on a smooth
complex projective manifold and outline the closure argument proving algebraicity
once the companion construction theorems are granted.
The core ingredient is a \emph{realization/microstructure} scheme producing
integral cycles whose calibration defect tends to~$0$ and whose masses converge
to the cohomological lower bound.  Any calibrated limit is a positive combination
of complex subvarieties~\cite{King71,HL82}, hence algebraic on projective
manifolds~\cite{Chow49,Serre56}.  A signed decomposition reduces the general
case to the cone-positive case, and Hard Lefschetz~\cite{Voisin02,GH78} treats
the full $p$-range.

\subsection*{Problem}

Let $X$ be a smooth projective complex variety of complex dimension $n$,
equipped with a K\"ahler form $\omega$ in the Chern class of a fixed ample line
bundle $L\to X$.  Fix an integer $1 \leq p \leq n$ and a rational Hodge class
\[
\gamma \;\in\; H^{2p}(X,\Q) \cap H^{p,p}(X).
\]
The Hodge conjecture asks whether $\gamma$ is represented by an algebraic cycle
of codimension~$p$.

\subsection*{Route}

Set $\psi:=\omega^{n-p}/(n-p)!$, the K\"ahler calibration of complex
$(n{-}p)$-planes, and $k:=2n-2p$.  For an integral $k$-cycle~$T$, the
\emph{calibration defect} is
\[
\Defcal(T):=\Mass(T)-\int_T\psi\;\ge\;0.
\]
The proof has three stages.

\paragraph{1. Signed decomposition.}
Every rational Hodge class $\gamma$ decomposes as
$\gamma=\gamma^+-\gamma^-$ where $\gamma^+$ is cone-positive and
$\gamma^-=N[\omega^p]$ is algebraic (Lemma~\ref{lem:signed-decomp}).

\paragraph{2. Cone-positive realization (SYR).}
For a cone-positive class $\gamma^+$ with smooth closed cone-valued
representative $\beta$, the companion papers construct integral
$k$-cycles $T_j$ with $\partial T_j=0$, $[T_j]=\PD(m[\gamma^+])$,
and $\Defcal(T_j)\to 0$ (Theorem~\ref{thm:automatic-syr}).

\paragraph{3. Calibrated limit and algebraicity.}
A subsequential flat limit of $\{T_j\}$ is $\psi$-calibrated, hence a
holomorphic chain by Harvey--Lawson/King, and algebraic by Chow/GAGA
(Theorem~\ref{thm:realization-from-almost}).

\subsection*{Organization}

Section~\ref{sec:interface} states the interface theorems from the companion
papers (without proof).
Section~\ref{sec:closure} proves the closure chain: almost-calibrated sequences
have calibrated limits; SYR-realizable forms yield algebraic cycles.
Section~\ref{sec:syr} states the Automatic SYR theorem (proved by
combining the companion-paper interfaces).
Section~\ref{sec:signed} proves the signed decomposition.
Section~\ref{sec:main} states and proves the main theorem.

\paragraph{Series convention.}
All results in Sections~\ref{sec:closure}--\ref{sec:main} are proved from the
interface propositions in Section~\ref{sec:interface}. Those interface results
are established in the companion preprints
``Bergman-scale holomorphic complete intersections with prescribed tangent
directions and template-controlled face traces'' (Paper~I, Theorem~1.1)~\cite{Washburn26I},
``Coherent discretization of strongly positive forms by vertex-prefix face
balancing'' (Paper~II, Theorem~1.1)~\cite{Washburn26II}, and
``Fixed-class integral cycles modulo torsion from coherent calibrated packets via period-first
rounding'' (Paper~III, Theorem~1.1)~\cite{Washburn26III}. The dependence on those inputs is intended to
be explicit rather than implicit.


%%% ============================================================
\section{Interface theorems from companion papers}\label{sec:interface}
%%% ============================================================

The following three propositions are the imported outputs of the companion series.
They are used in this paper only as explicit interface inputs.

\begin{proposition}[H1 interface: local holomorphic manufacturing]
\label{prop:h1-package}
For each mesh cell $Q$ and each direction family prescribed by the local
Carath\'eodory data of $\beta$ on $Q$, there exists a $\psi$-calibrated
sheet-sum integral current $S_Q$ satisfying $\Mass(S_Q)=\langle S_Q,\psi\rangle$
with quantitative slope and budget control, together with template-controlled
face traces in flat norm.
\end{proposition}

\begin{proof}[Source in the series]
This is exactly Theorem~1.1 of Paper~I,
``Bergman-scale holomorphic complete intersections with prescribed tangent
directions and template-controlled face traces''~\cite{Washburn26I}.
\end{proof}

\begin{proposition}[H2 interface: assembly packet existence]
\label{prop:h2-package}
For every sufficiently small mesh size $h>0$, there exists an assembly packet
at scale $h$ for $\beta$ satisfying the three properties needed downstream:

\begin{enumerate}
\item[\textnormal{(A1)}] cube-wise mass and barycenter control,
\item[\textnormal{(A2)}] global coherence across labels, and
\item[\textnormal{(A3)}] a marginal period bound.
\end{enumerate}
\end{proposition}

\begin{proof}[Source in the series]
This is exactly Theorem~1.1 of Paper~II,
``Coherent discretization of strongly positive forms by vertex-prefix face
balancing''~\cite{Washburn26II}.
\end{proof}

\begin{proposition}[H3 interface: exact-class gluing and almost-calibration]
\label{prop:h3-package}
Assume an assembly packet satisfying \textnormal{(A1)}--\textnormal{(A3)} as in
Proposition~\ref{prop:h2-package}. Then, for all sufficiently small mesh sizes,
one can choose the final marginal activations so as to obtain a corrected
integral current $T_h$ with:
\begin{enumerate}
\item[\textnormal{(i)}] $\partial T_h=0$,
\item[\textnormal{(ii)}] $[T_h]=\PD(m[\gamma])$ in $H_k(X,\Z)/\mathrm{tors}$,
\item[\textnormal{(iii)}] $0\le \Defcal(T_h)\to 0$, and hence
\item[\textnormal{(iv)}] $\Mass(T_h)\to c_0$.
\end{enumerate}
\end{proposition}

\begin{proof}[Source in the series]
This is exactly Theorem~1.1 of Paper~III,
``Fixed-class integral cycles modulo torsion from coherent calibrated packets via period-first
rounding''~\cite{Washburn26III}.
\end{proof}


%%% ============================================================
\section{Closure chain}\label{sec:closure}
%%% ============================================================

\subsection*{Spine theorem}

\begin{theorem}[Quantitative almost-mass-minimizing cycles]
\label{thm:spine-quantitative}
Let $(X,\omega)$ be a smooth projective K\"ahler manifold of complex
dimension $n$, fix $1\le p\le n$, and set $\psi=\omega^{n-p}/(n-p)!$ and
$k:=2n-2p$.  Let $[\gamma]\in H^{2p}(X,\Q)\cap H^{p,p}(X)$ admit a smooth
closed cone-valued representative $\beta$.  Choose $m\ge 1$ so that
$m[\gamma]\in H^{2p}(X,\Z)$, and set
\[
A:=\PD(m[\gamma])\in H_k(X,\Z),
\qquad
c_0:=\langle A,[\psi]\rangle=m\int_X \beta\wedge\psi.
\]
For each mesh scale $h_j\to 0$, applying Propositions~\ref{prop:h1-package},
\ref{prop:h2-package}, and~\ref{prop:h3-package} produces integral $k$-cycles $T_j$ with
$\partial T_j=0$, $[T_j]=A$, and
\[
0\;\le\;\Mass(T_j)-c_0\;\le\;2\,\Mass(G_j),
\]
where $G_j$ is the gluing correction with $\Mass(G_j)\to 0$.
In particular $\Mass(T_j)\to c_0$, i.e.\ $\Defcal(T_j)\to 0$.
\end{theorem}

\begin{proof}
Since $\psi$ is closed and $[T_j]=A$, the pairing $\langle T_j,\psi\rangle
=\langle A,[\psi]\rangle=c_0$ is topological.  The calibration inequality
$\Mass(T_j)\ge c_0$ gives the lower bound.  Writing $S_j=T_j+G_j$ where
$S_j$ is $\psi$-calibrated (so $\Mass(S_j)=\langle S_j,\psi\rangle$), the
triangle inequality yields
\[
\Mass(T_j)\le \Mass(S_j)+\Mass(G_j)
=\langle T_j+G_j,\psi\rangle+\Mass(G_j)
\le c_0+2\,\Mass(G_j).
\qedhere
\]
\end{proof}


\subsection*{Microstructure/gluing estimate}

\begin{proposition}[Microstructure/gluing estimate]\label{prop:glue-gap}
Let $T^{\mathrm{raw}}$ be the integral $k$-current built from the
microstructure pieces on a mesh of size $h$, $R:=\partial T^{\mathrm{raw}}$,
and $\delta:=\mathcal F(R)$.  Then there exists an integral $k$-current
$R_{\mathrm{glue}}$ with
\[
\partial R_{\mathrm{glue}}=-R,
\qquad
\Mass(R_{\mathrm{glue}})\le \delta + C_X\,\delta^{\frac{k}{k-1}},
\]
where $C_X$ depends only on $X$ and $k$.
\end{proposition}

\begin{proof}
By the flat-norm definition, decompose $R=R_0+\partial Q$ with
$\Mass(R_0)+\Mass(Q)\le\delta+\eta$.
Then $R_0$ is an integral $(k{-}1)$-cycle bounding in $X$, so the
Federer--Fleming filling inequality~\cite{FF60,Fed69} gives $Q_0$ with
$\partial Q_0=R_0$ and $\Mass(Q_0)\le C_X\Mass(R_0)^{k/(k-1)}$.
Setting $R_{\mathrm{glue}}:=-(Q+Q_0)$ and letting $\eta\to 0$ yields
the bound.
\end{proof}


\subsection*{Cohomology matching}

\begin{proposition}[Integral cohomology constraints]\label{prop:cohomology-match}
Fix smooth closed $k$-forms $\Theta_1,\dots,\Theta_b$ representing an integral
basis of $H^k(X,\Z)/\mathrm{tors}$, and write
\[
I_\ell:=\int_X \beta\wedge\Theta_\ell.
\]
Let $S$ be the raw integral current obtained after integer rounding, and let
$U_\eps$ be an integral current with $\partial U_\eps=\partial S$ and
$\Mass(U_\eps)\to 0$. Set $T_\eps:=S-U_\eps$.
By refining the mesh and choosing the integer sheet counts $N_{Q,j}$ via
fixed-dimensional discrepancy rounding~\cite{BaranyGrinberg81}, one achieves
\[
\int_{T_\eps}\Theta_\ell = m\,I_\ell
\qquad (\ell=1,\dots,b),
\]
so that $[T_\eps]=\PD(m[\gamma])$ in $H_k(X,\Z)/\mathrm{tors}$.
\end{proposition}

\begin{proof}
Cube-wise mass and barycenter control force the fractional bookkeeping
current to approximate the target period vector within $\frac18$.
Discrepancy rounding adds at most $\frac18$ more.
The gluing correction contributes $<\frac14$ to each period.
Since closed integral cycles pair integrally with integral classes, the
total error $<1$ forces exact equality.
\end{proof}


\subsection*{Almost-calibration}

\begin{proposition}[Almost-calibration]\label{prop:almost-calibration}
Let $S$ be a $\psi$-calibrated sheet sum and $U_\eps$ an integral current
with $\partial U_\eps=\partial S$ and $\Mass(U_\eps)\to 0$.  Set
$T_\eps:=S-U_\eps$.  Then:
\begin{enumerate}
\item[\textnormal{(i)}] $\displaystyle\int_{T_\eps}\psi
  =\langle\PD(m[\gamma]),[\psi]\rangle=c_0$ for all~$\eps$.
\item[\textnormal{(ii)}] $0\le\Defcal(T_\eps)\le 2\,\Mass(U_\eps)\to 0$.
\item[\textnormal{(iii)}] $\Mass(T_\eps)\to c_0$.
\end{enumerate}
\end{proposition}

\begin{proof}
Item (i): $\psi$ is closed and $[T_\eps]$ is fixed, so the pairing is
topological.
For (ii), $\Mass(S)=\int_S\psi$ and the comass bound give
\[
\Defcal(T_\eps)
=\Mass(T_\eps)-\int_{T_\eps}\psi
\le(\Mass(S)+\Mass(U_\eps))-(\int_S\psi-|\int_{U_\eps}\psi|)
\le 2\Mass(U_\eps).
\]
Item (iii) combines (i) and (ii).
\end{proof}


\subsection*{Calibrated limits are holomorphic}

\begin{definition}[Stationary Young-measure realizability (SYR)]
\label{def:syr}
A cone-valued smooth closed $(p,p)$-form $\beta$ representing $[\gamma]$ is
\emph{SYR-realizable} if there exist integral $(2n{-}2p)$-cycles $T_k$ with
$\partial T_k=0$, $[T_k]=\PD(m[\gamma])$ for a fixed $m\ge 1$, and
$\Defcal(T_k)\to 0$.
\end{definition}

\begin{theorem}[Realization from almost-calibrated sequences]
\label{thm:realization-from-almost}
Assume $\beta$ is SYR-realizable.  Then there exists an integral
$(2n{-}2p)$-cycle $T$ with $\partial T=0$, $[T]=\PD(m[\gamma])$, and
$\Mass(T)=\langle T,\psi\rangle=c_0$.

In particular, $T$ is $\psi$-calibrated.  By Harvey--Lawson~\cite{HL82},
$T$ is a positive $d$-closed integral current of bidimension $(p,p)$.
By King's theorem~\cite{King71}, $T$ is a holomorphic chain
$T=\sum_j m_j[V_j]$ with $V_j$ irreducible analytic subvarieties.
If $X$ is projective, each $V_j$ is algebraic by Chow/GAGA~\cite{Chow49,Serre56}.
\end{theorem}

\begin{proof}
Since $[T_k]=\PD(m[\gamma])$ and $\psi$ is closed,
$\langle T_k,\psi\rangle=c_0$ for all $k$.
The defect hypothesis gives $\Mass(T_k)\to c_0$, so
$\sup_k\Mass(T_k)<\infty$.
By Federer--Fleming compactness~\cite{Fed69}, passing to a subsequence,
$T_k\to T$ in flat norm with $T$ an integral cycle and $[T]=\PD(m[\gamma])$.
Flat convergence gives $\langle T,\psi\rangle=c_0$; lower semicontinuity
gives $\Mass(T)\le c_0$; the comass inequality gives
$\langle T,\psi\rangle\le\Mass(T)$.
Combining: $\Mass(T)=c_0=\langle T,\psi\rangle$, so $T$ is $\psi$-calibrated.
The holomorphic-chain conclusion and algebraicity follow from
Harvey--Lawson~\cite{HL82}, King~\cite{King71}, and Chow/GAGA~\cite{Chow49,Serre56,Hartshorne77,GH78}.
\end{proof}

\begin{theorem}[Calibrated realization under SYR]\label{thm:syr}
Assume $\beta$ is SYR-realizable.  Then the conclusion of
Theorem~\ref{thm:realization-from-almost} holds: there exists a
$\psi$-calibrated holomorphic chain $T=\sum_j m_j[V_j]$ representing
$\PD(m[\gamma])$, and $[\gamma]$ is algebraic if $X$ is projective.
\end{theorem}

\begin{proof}
Apply Theorem~\ref{thm:realization-from-almost} to the SYR sequence.
\end{proof}


%%% ============================================================
\section{Automatic SYR}\label{sec:syr}
%%% ============================================================

\begin{theorem}[Automatic SYR for cone-valued forms]\label{thm:automatic-syr}
Assume Propositions~\ref{prop:h1-package}, \ref{prop:h2-package}, and~\ref{prop:h3-package}.
Let $(X,\omega)$ be a smooth complex projective manifold of complex
dimension $n$, and let $1\le p\le n/2$.  (For $p>n/2$, reduce to the
complementary degree $n-p$ by Hard Lefschetz~\cite{Voisin02,GH78}.)
Let $\beta$ be a smooth closed cone-valued $(p,p)$-form representing
$[\gamma]\in H^{p,p}(X;\Q)$.

Then $\beta$ is SYR-realizable: there exist integral $(2n{-}2p)$-cycles
$T_k$ with $\partial T_k=0$,
$[T_k]=\PD(m[\gamma])$, and $\Defcal(T_k)\to 0$.
In particular, by Theorem~\ref{thm:syr}, $[\gamma]$ is algebraic.
\end{theorem}

\begin{proof}
Fix a mesh parameter $\eps>0$.
By Proposition~\ref{prop:h2-package}, for each sufficiently small scale
$\eps>0$ there exists an assembly packet satisfying \textnormal{(A1)}--\textnormal{(A3)}.
Applying Proposition~\ref{prop:h3-package} to that packet yields closed
integral cycles $T_\eps$ with $[T_\eps]=\PD(m[\gamma])$ and
$\Defcal(T_\eps)\to 0$ as $\eps\to 0$.
Choosing $\eps_k\downarrow 0$ gives the SYR sequence.
Theorem~\ref{thm:syr} concludes algebraicity.
\end{proof}


%%% ============================================================
\section{Signed decomposition}\label{sec:signed}
%%% ============================================================

\begin{lemma}[Strict interior positivity of the K\"ahler power]
\label{lem:kahler-positive}
The $(p,p)$-form $\omega^p$ is strictly positive in the calibrated cone:
$\omega^p(x)\in\mathrm{int}\,K_p(x)$ for all $x\in X$.
There exists a uniform radius $r_0>0$ such that
$B(\omega^p(x),r_0)\subset K_p(x)$ for every $x$.
\end{lemma}

\begin{proof}
In a unitary frame at $x$, $\omega^p_x$ is a strictly (strongly) positive
$(p,p)$-form.
Compactness of $X$ and continuity of the cone give uniform $r_0>0$.
\end{proof}

\begin{lemma}[Signed decomposition]\label{lem:signed-decomp}
Let $\gamma\in H^{2p}(X,\Q)\cap H^{p,p}(X)$ be any rational Hodge class.
Then there exist cone-positive rational Hodge classes $\gamma^+$ and
$\gamma^-$ with $\gamma=\gamma^+-\gamma^-$, where
$\gamma^-=N[\omega^p]$ for some $N\in\Q_{>0}$.
\end{lemma}

\begin{proof}
Let $\alpha$ represent $\gamma$, set $M:=\sup_X\|\alpha\|$, and choose
$N\in\Q_{>0}$ with $N>M/r_0$.
Then $\alpha(x)+N\omega^p(x)\in K_p(x)$ for all $x$, so
$\gamma^+:=\gamma+N[\omega^p]$ is cone-positive and
$\gamma^-:=N[\omega^p]$ is algebraic.
\end{proof}

\begin{lemma}[$\omega^p$ is algebraic]\label{lem:gamma-minus-alg}
The class $[\omega^p]=c_1(L)^p$ lies in the $\Q$-span of algebraic
cycle classes.  For $\mhol\gg 0$, a smooth complete intersection
$Z=D_1\cap\cdots\cap D_p$ with $D_i\in|L^{\otimes\mhol}|$ satisfies
$\PD([Z])=\mhol^p[\omega^p]$.
\end{lemma}

\begin{proof}
Very ampleness of $L^{\otimes\mhol}$ and Bertini's
theorem~\cite{Hartshorne77} give a smooth $Z$ with the stated class.
\end{proof}


%%% ============================================================
\section{Main theorem}\label{sec:main}
%%% ============================================================

\begin{theorem}[Cone-positive classes are algebraic]
\label{thm:effective-algebraic}
Assume Propositions~\ref{prop:h1-package}, \ref{prop:h2-package}, and~\ref{prop:h3-package}.
Let $\gamma^+\in H^{2p}(X,\Q)\cap H^{p,p}(X)$ be a cone-positive
rational Hodge class on a smooth complex projective manifold, with
$p\le n/2$.  Then $\gamma^+$ is algebraic.
\end{theorem}

\begin{proof}
By Theorem~\ref{thm:automatic-syr}, $\gamma^+$ is SYR-realizable.
By Theorem~\ref{thm:syr}, there exists a $\psi$-calibrated holomorphic
chain representing $\PD(m[\gamma^+])$.
Since $X$ is projective, Chow/GAGA gives algebraicity.
\end{proof}

\begin{conjecture}[Hodge Conjecture]\label{conj:hodge-conjecture}
Let $X$ be a smooth complex projective manifold.  Every rational
cohomology class $\gamma\in H^{2p}(X,\Q)\cap H^{p,p}(X)$ is
represented by an algebraic cycle.
\end{conjecture}

\begin{theorem}[Hodge Conjecture for rational $(p,p)$ classes]
\label{thm:main-hodge}
Assume Propositions~\ref{prop:h1-package}, \ref{prop:h2-package}, and~\ref{prop:h3-package}.
Let $X$ be a smooth complex projective manifold.  Then every rational
Hodge class $\gamma\in H^{2p}(X,\Q)\cap H^{p,p}(X)$ is algebraic.
\end{theorem}

\begin{proof}
By Hard Lefschetz~\cite{Voisin02,GH78} it suffices to treat $p\le n/2$.
Apply Lemma~\ref{lem:signed-decomp}: $\gamma=\gamma^+-\gamma^-$ with
$\gamma^-=N[\omega^p]$ algebraic (Lemma~\ref{lem:gamma-minus-alg}) and
$\gamma^+$ cone-positive.
By Theorem~\ref{thm:effective-algebraic}, $\gamma^+$ is algebraic.
Since algebraic classes form a $\Q$-subspace,
$\gamma=\gamma^+-\gamma^-$ is algebraic.
\end{proof}

\begin{corollary}[Full Hodge conjecture]\label{cor:full-hodge}
Every rational $(p,p)$ class on a smooth complex projective manifold is
represented by an algebraic cycle.
\end{corollary}

\begin{proof}
Immediate from Theorem~\ref{thm:main-hodge}.
\end{proof}

\begin{remark}[Logical structure]
\label{rem:hodge-final-remark}
Theorem~\ref{thm:main-hodge} combines three inputs:
(i) signed decomposition (Lemma~\ref{lem:signed-decomp});
(ii) SYR realization of cone-positive classes
(Theorem~\ref{thm:automatic-syr}, proved using the companion-paper
interfaces);
(iii) the calibrated-limit closure principle
(Theorem~\ref{thm:realization-from-almost}) with algebraicity via
Chow/GAGA.
\end{remark}


%%% ============================================================
\begin{thebibliography}{99}

\bibitem{Washburn26I}
J.~Washburn.
\newblock {\em Bergman-scale holomorphic complete intersections with prescribed
  tangent directions and face traces}.
\newblock Preprint, 2026.

\bibitem{Washburn26II}
J.~Washburn.
\newblock {\em Coherent discretization of strongly positive forms by
  vertex-prefix face balancing}.
\newblock Preprint, 2026.

\bibitem{Washburn26III}
J.~Washburn.
\newblock {\em Exact-class integral cycles from coherent calibrated packets via
  period-first rounding}.
\newblock Preprint, 2026.

\bibitem{BaranyGrinberg81}
I.~B\'ar\'any and V.~S.~Grinberg.
\newblock On some combinatorial questions in finite-dimensional spaces.
\newblock {\em Linear Algebra Appl.}, 41:1--9, 1981.

\bibitem{Chow49}
W.-L. Chow.
\newblock On compact complex analytic varieties.
\newblock {\em Amer.\ J.\ Math.}, 71:893--914, 1949.

\bibitem{FF60}
H.~Federer and W.~H.~Fleming.
\newblock Normal and integral currents.
\newblock {\em Annals of Mathematics}, 72(3):458--520, 1960.

\bibitem{Fed69}
H.~Federer.
\newblock {\em Geometric Measure Theory}.
\newblock Springer, 1969.

\bibitem{GH78}
P.~Griffiths and J.~Harris.
\newblock {\em Principles of Algebraic Geometry}.
\newblock Wiley-Interscience, 1978.

\bibitem{Hartshorne77}
R.~Hartshorne.
\newblock {\em Algebraic Geometry}.
\newblock Graduate Texts in Mathematics 52, Springer, 1977.

\bibitem{HL82}
R.~Harvey and H.~B.~Lawson, Jr.
\newblock Calibrated geometries.
\newblock {\em Acta Mathematica}, 148:47--157, 1982.

\bibitem{King71}
J.~R.~King.
\newblock The currents defined by analytic varieties.
\newblock {\em Acta Mathematica}, 127:185--220, 1971.

\bibitem{Serre56}
J.-P.~Serre.
\newblock G\'eom\'etrie alg\'ebrique et g\'eom\'etrie analytique (GAGA).
\newblock {\em Annales de l'Institut Fourier}, 6:1--42, 1956.

\bibitem{Voisin02}
C.~Voisin.
\newblock {\em Hodge Theory and Complex Algebraic Geometry I, II}.
\newblock Cambridge University Press, 2002--2003.

\end{thebibliography}

\end{document}
